\documentclass{article}
\usepackage[top=2cm,bottom=2cm,left=3cm,right=3cm]{geometry}
\usepackage{enumitem}
\usepackage{xcolor}
\usepackage{minted}
\usepackage[most]{tcolorbox}
\tcbuselibrary{minted,breakable}
\usepackage{fontspec}
\usepackage{polyglossia}
\setdefaultlanguage{ukrainian}
\setotherlanguages{english}
\usepackage{Alegreya}
\usepackage{FiraMono}
\newfontfamily\cyrillicfont{Alegreya}
\newfontfamily\cyrillicfontsf{Alegreya Sans}
\newfontfamily\cyrillicfonttt{DejaVu Sans Mono}
\usepackage{sciffi-python}
\usepackage{sciffi-memo}

\usemintedstyle{emacs}

\newtcolorbox{CodeBox}[1][]{
    breakable,
    enhanced,
    colback=gray!10,
    colframe=gray!70,
    fonttitle=\bfseries,
    title=#1,
    sharp corners,
    boxrule=0.8pt,
    top=2mm,
    bottom=2mm,
    left=2mm,
    right=2mm,
    fontupper=\footnotesize,
}

\newtcolorbox{ResultBox}[1][]{
    breakable,
    enhanced,
    colback=gray!5,
    colframe=gray!50,
    fonttitle=\bfseries,
    title=#1,
    sharp corners,
    boxrule=0.6pt,
    top=1mm,
    bottom=1mm,
    left=2mm,
    right=2mm,
    fontupper=\footnotesize,
}

\begin{document}
    \begin{titlepage}
        \begin{center}
            {\Large
                МІНІСТЕРСТВО ОСВІТИ ТА НАУКИ УКРАЇНИ

                НАЦІОНАЛЬНИЙ ТЕХНІЧНИЙ УНІВЕРСИТЕТ УКРАЇНИ
                «КИЇВСЬКИЙ ПОЛІТЕХНІЧНИЙ ІНСТИТУТ ІМЕНІ ІГОРЯ СІКОРСЬКОГО»

                ФАКУЛЬТЕТ ІНФОРМАТИКИ ТА ОБЧИСЛЮВАЛЬНОЇ ТЕХНІКИ

                КАФЕДРА ІНФОРМАЦІЙНИХ СИСТЕМ ТА ТЕХНОЛОГІЙ
            }


            \vspace{60mm}
            {\large
                \textbf{ЗВІТ}

                \vspace{5mm}

                \textbf{До лабораторної роботи 4}

                \vspace{5mm}

                з дисципліни «Технології Розробки Програмного Забезпечення»
            }

            На тему «Музичний Магазин»

            варіант \textnumero 14 (\textnumero 4)
        \end{center}

        \vfill
        \hfill
        \begin{minipage}[t]{0.35\textwidth}
            \textbf{Виконав:}
        \end{minipage}

        \vfill

        \begin{center}
            {\bf
                Київ

                КПІ Ім. Ігоря Сікорського

                2025
            }
        \end{center}
    \end{titlepage}
    \newpage

    Тема: підзапити до даних БД.

    Мета: вивчити команди формування підзапитів до бази даних.

    Навички що будуть здобуті: вміння формувати складні запити до бази даних, що складаються з одного або декількох підзапитів.

    \section*{Хід виконання:}
    \begin{sciffi}{python}
        from pathlib import Path
        from subprocess import run, PIPE
        from collections.abc import Sequence
        from contextlib import contextmanager
        
        @contextmanager
        def env(name: str, *args: str):
            print(f"\\begin{{{name}}}" + "".join(args))
            try:
                yield
            finally:
                print(f"\\end{{{name}}}")

        def sh(*args: str | Path) -> str:
            return run(
                tuple(map(str, args)),
                stdout=PIPE,
                stderr=PIPE,
                text=True,
            ).stdout

        scripts = Path("./scripts/")
        # Find all .sql files and sort them numerically by their filename
        sql_files = sorted(scripts.glob("*.sql"), key=lambda p: int(p.stem))

        flow: Sequence[tuple[Path, Sequence[str]]] = [
            (script, ("--no-echo",)) for script in sql_files
        ]

        for script, args in flow:
            source = script.read_text(encoding="utf-8")
            output = sh("tsqlexec", *args, script)

            print(f"\\subsection*{{Запит з файлу: {script.name}}}")
            with env("CodeBox", "[SQL Source]"):
                with env("minted", "[breaklines,linenos]", "{sql}"):
                    print(source, end="")
            print()

            if not output.strip():
                continue

            print("Result:")
            with env("ResultBox", "[Output]"):
                with env("minted", "[breaklines]", "{text}"):
                    print(output, end="")
            print("\\vspace{5mm}")

    \end{sciffi}

    \newpage
    \section*{Контрольні питання}
    \begin{enumerate}[label=\arabic*., leftmargin=*]
        \item \textbf{Що таке підзапит?}
        
        Підзапит (subquery) — це запит \texttt{SELECT}, вкладений всередину іншого SQL-запиту (головного запиту). Результат підзапиту використовується головним запитом для виконання певної операції (фільтрації, порівняння, як джерело даних тощо).

        \item \textbf{В яких командах SQL можуть бути використані підзапити?}
        
        Підзапити можуть використовуватися в основних командах DML (Data Manipulation Language):
        \begin{itemize}
            \item \texttt{SELECT} (у пропозиціях \texttt{SELECT}, \texttt{FROM}, \texttt{WHERE}, \texttt{HAVING})
            \item \texttt{INSERT} (для надання значень для вставки)
            \item \texttt{UPDATE} (у пропозиціях \texttt{SET} або \texttt{WHERE})
            \item \texttt{DELETE} (у пропозиції \texttt{WHERE})
        \end{itemize}

        \item \textbf{Який порядок виконання підзапита?}
        
        Зазвичай, підзапит виконується \textbf{першим}, а його результат передається зовнішньому (головному) запиту для подальшої обробки. Це стосується некорельованих підзапитів. У випадку корельованих підзапитів, внутрішній запит виконується багаторазово — для кожного рядка, що обробляється зовнішнім запитом.

        \item \textbf{В яких випадках підзапити можуть передати декілька рядків в головну пропозицію?}
        
        Підзапити повертають декілька рядків, коли вони використовуються з операторами, призначеними для роботи з наборами даних, такими як \texttt{IN}, \texttt{NOT IN}, \texttt{ANY}, \texttt{ALL}, \texttt{EXISTS}. Також підзапит у пропозиції \texttt{FROM} (похідна таблиця) за своєю природою повертає набір рядків.

        \item \textbf{Де розміщується підзапит в логічному виразі пропозиції основного запита: до оператора порівняння, після оператора порівняння?}
        
        Підзапит завжди розміщується \textbf{після} оператора порівняння.
        \begin{CodeBox}[Приклад]
        \begin{minted}[breaklines]{sql}
WHERE column_name > (SELECT MIN(column) FROM ...)
        \end{minted}
        \end{CodeBox}

        \item \textbf{Які пропозиції допустимі для основного запита не можуть бути використані в підзапиті?}
        
        Найчастіше в підзапиті не можна використовувати пропозицію \texttt{ORDER BY}. Сортування результату підзапиту не має сенсу, оскільки кінцевий порядок рядків визначає головний запит. Винятком є випадки, коли \texttt{ORDER BY} використовується разом з командами обмеження кількості рядків (\texttt{TOP}, \texttt{LIMIT}, \texttt{FETCH}).

        \item \textbf{Які існують типи підзапитів?}
        
        Підзапити можна класифікувати за кількома критеріями:
        \begin{itemize}
            \item \textbf{За результатом, що повертається:}
            \begin{itemize}
                \item \textit{Скалярний (однорядковий):} повертає один рядок і один стовпець.
                \item \textit{Багаторядковий:} повертає декілька рядків, але один стовпець.
                \item \textit{Табличний (багатостовпцевий):} повертає декілька рядків і декілька стовпців (використовується у \texttt{FROM}).
            \end{itemize}
            \item \textbf{За залежністю від зовнішнього запиту:}
            \begin{itemize}
                \item \textit{Некорельований:} може бути виконаний незалежно від зовнішнього запиту.
                \item \textit{Корельований:} залежить від значень у рядках зовнішнього запиту.
            \end{itemize}
        \end{itemize}

        \item \textbf{В яких пропозиціях SELECT можуть бути використані підзапити?}
        
        У команді \texttt{SELECT} підзапити можуть бути використані в таких пропозиціях:
        \begin{itemize}
            \item \texttt{SELECT list} (скалярний підзапит для обчислення значення стовпця).
            \item \texttt{FROM} (табличний підзапит, що створює похідну таблицю).
            \item \texttt{WHERE} (для фільтрації рядків).
            \item \texttt{HAVING} (для фільтрації груп).
        \end{itemize}

        \item \textbf{Які пропозиції SELECT неможна використовувати в підзапитах?}
        
        Як згадувалося раніше, пропозицію \texttt{ORDER BY} зазвичай не можна використовувати в підзапитах. Також деякі СУБД можуть накладати обмеження на використання \texttt{INTO} (для створення нової таблиці) всередині підзапиту.

        \item \textbf{Де повинен знаходитись підзапит?}
        
        Синтаксично підзапит завжди повинен бути взятий у \textbf{круглі дужки} \texttt{()}.

        \item \textbf{Які оператори порівняння працюють с декількома рядками?}
        
        Оператори, що працюють з набором значень (декількома рядками), які повертає підзапит:
        \begin{itemize}
            \item \texttt{IN} — дорівнює будь-якому значенню зі списку.
            \item \texttt{NOT IN} — не дорівнює жодному значенню зі списку.
            \item \texttt{ANY} / \texttt{SOME} — порівнює значення з кожним значенням зі списку (використовується з \texttt{=}, \texttt{>}, \texttt{<} і т.д.).
            \item \texttt{ALL} — порівнює значення з усіма значеннями зі списку.
        \end{itemize}

        \item \textbf{Які оператори порівняння працюють тільки з одним рядком?}
        
        Стандартні оператори порівняння, які очікують, що підзапит поверне одне скалярне значення:
        \texttt{=}, \texttt{<>} (\texttt{!=}), \texttt{>}, \texttt{<}, \texttt{>=}, \texttt{<=}.

        \item \textbf{Які підзапити називаються однорядковими?}
        
        Однорядковими (скалярними) називаються підзапити, які гарантовано повертають не більше одного рядка та одного стовпця.

        \item \textbf{Які підзапити називаються квантифікованими?}
        
        Квантифікованими називаються підзапити, що використовуються з операторами-кванторами \texttt{ANY}, \texttt{SOME} або \texttt{ALL}. Ці оператори дозволяють порівняти значення з набором значень, які повертає підзапит.

        \item \textbf{В чому різниця підзапитів у FROM та WHERE?}
        
        Основна різниця полягає в їхній ролі:
        \begin{itemize}
            \item \textbf{Підзапит у \texttt{WHERE} (або \texttt{HAVING}):} Використовується для \textbf{фільтрації}. Він повертає значення або набір значень, які використовуються в умові для відбору рядків з таблиці, вказаної у \texttt{FROM} основного запиту.
            \item \textbf{Підзапит у \texttt{FROM}:} Використовується як \textbf{джерело даних}. Він створює тимчасову, віртуальну таблицю (похідну таблицю), з якої зовнішній запит вибирає дані. До цієї таблиці можна застосовувати \texttt{JOIN}, \texttt{WHERE} та інші операції, як до звичайної таблиці.
        \end{itemize}
    \end{enumerate}
\end{document}
